\section{Simple Test of Linear Divergence}\label{sec:test}

To study the renormalization properties of the quasi-PDF operator, we need to calculate matrix elements at several different lattice spacings. To ensure the data
to be useful for a refined high-precision analysis, we start by testing whether they approximately show the linear divergence  predicted by perturbation theory. In particular, one needs to show that the $z$-dependence of the linearly divergent term is linear. 

To achieve this, we first extract the factor $e^{-\delta \bar{m} z}$ from the bare matrix element to test the $1/a$ dependence in $\delta \bar{m}$ in Eq.~(\ref{eq:delta_m}). 
Based on Eqs.~(\ref{eq:Reoperator}),~(\ref{eq:delta_m}), and~(\ref{eq:alpha_s}), if we take the natural log on a bare matrix element $\mathcal{M}$, we have,
\bea\label{eq:logM_s}
\ln \mathcal{M}(z,a) = \frac{e(z)}{a \ln[a \Lambda_{\rm QCD}]} + g(z),
\eea
where the first term is linearly divergent and $g(z)$ is the residual. We have ignored the logarithmic divergence and discretization error here because these terms are much smaller and thus have little influence on this simple test of the linear divergence.

\begin{figure}[tbp]
\centering
\includegraphics[width=8cm]{images/fitting_RIMOM_overlap.pdf}
\includegraphics[width=8cm]{images/ez_RIMOM_overlap.pdf}
\caption{
Upper Panel: Using Eq.~(\ref{eq:logM_s}) to fit the bare $O_{\gamma_t}(z)$ matrix element  in the off-shell quark state for overlap action without HYP smearing. Blue points are interpolated data and colorful curves are fitted curves for each $z$. $\Lambda_{\rm QCD}$ is fixed at 0.2 GeV.
Lower Panel: $e(z)$ with respect to $z$. Blue points are the fitted parameters $e(z)$ from the upper panel. The blue curve is the no-bias linear fit for $e(z)$.}
\label{fig:testlindiv_RIMOM_overlap}
\end{figure}

\begin{figure}[tbp]
\centering
\includegraphics[width=8cm]{images/fitting_RIMOM_clover.pdf}
\includegraphics[width=8cm]{images/ez_RIMOM_clover.pdf}
\caption{Same as Fig.~\ref{fig:testlindiv_RIMOM_overlap}, except for the clover valence quark.}
\label{fig:testlindiv_RIMOM_clover}
\end{figure}

We use Eq.~(\ref{eq:logM_s}) to fit the bare matrix elements as a function of $a$. We treat $e(z)$ and $g(z)$ as unknown functions of $z$ so that our fitting is performed at each value of $z$ at which they are treated as free parameters. We need to do linear interpolation on the data $\ln \mathcal{M}(z,a)$ with respect to $z$. Here are the steps we perform in detail:

1. We do the linear interpolation with neighbourhood two data points to obtain the central value and uncertainty at $z=0.06\times n$ fm $(n=3,4,5,...,16)$.

2. For each $z$, we fit the dependence on $a$, treating $e(z)$, $g(z)$ as fitted parameters. For this initial test, we fix $\Lambda_{\rm QCD}$ at 0.2 GeV, which is a reasonable guess to start with~\cite{Lepage:1992xa}.

3. We plot the fitted $e(z)$ with respect to $z$ to see if $e(z)$ has a linear $z$ dependence.

If we obtain a reasonable fit in step 2, we then see approximately the $1/a$ dependence of the divergences. In step 3, we can test the linear $z$ dependence. The fitting results for the bare $O_{\gamma_t}(z)$ matrix element in the off-shell quark state for valence overlap and clover actions without HYP smearing are shown in Fig.~\ref{fig:testlindiv_RIMOM_overlap} and Fig.~\ref{fig:testlindiv_RIMOM_clover}, respectively. There
are eight different lattice spacings. The range of $z$
is taken from 0.18~fm to 0.96~fm and we analyse 14 different $z$ values in this range.

The $1/a$ dependence in both cases is fitted very well, although due to the high precision of the data, chi-square
is large. However, this is not our concern at this initial step. 
The $z$ dependence of the divergent coefficients
shows nicely the linear feature, approximately going through
zero at $z=0$. This is a strong indication that
the linear divergence follows roughly
the prediction of perturbation theory which gives
us confidence to perform more refined analyses in the next sections.
